\section{Discussion and lessons learned}
\label{sec:discussion}

The management evidence is the first result, not merely a precondition for the
retrieval experiments. Across the acquisition and organisation campaigns, HMRS
recorded 322 stored parts across five libraries, measured content-based
assignment separately from retrieval, and exposed state disagreement after a
rename and in selected lifecycle checks. The 79/21 arXiv availability result
and the 0.875/1.0/1.0 cross-lingual assignment result come from different
corpora and should be read as complementary evidence: the former identifies an
incomplete ingestion boundary, while the latter tests content-based placement.
The 33-part consistency sample and selected passing lifecycle checks show
that the management loop is exercisable, not that it is transactionally
complete. The failed task-status/update checks and stale-index rename are
important system results because they show why reconciliation and post-operation
probes are part of the design rather than optional diagnostics.

The access results then show how managed state is consumed. The frozen route-card
audit selected the gold collection in 72/72 full-card cases and 37/72 nominally
name-only cases, but its fallback and scripted nature prevent a causal
interpretation. In the same-question R3 content-layer follow-up, A0h, the tiered
configuration with a chunk-window union, had the highest observed exact-string
containment at 51/72 questions; PD, A0f, A0/A1b and A2 contained it in
45/72, 44/72, 39/72 and 25/72, respectively. A0h versus A0 was 12--0
($p=0.00049$, exploratory and unadjusted). These are containment proxies, not
semantic answers, and the configurations vary in more than navigation alone.
The frozen read-all scan found every string but generated no answers, while the
R3 read-all run was incomplete; this is why the mechanical upper bound is not a
current Agent baseline. The result supports the narrower claim that the
managed, tiered content configuration was effective on this corpus and question
set, not a universal hierarchy advantage.

The same-pool diagnostic illustrates how candidate-evidence ordering can matter
in this setting. On identical logged pools from the retired hybrid wrapper,
verifier-score ordering reproduced the Hit@5 decrease documented in
Section~\ref{sec:e6}. A separate score record contains 191 values above the
threshold; it does not establish why one candidate was absent from a ranking
replay, or the effect of filtering. This is a bounded warning about a
particular historical ordering, not a finding about every verifier or current
retrieval lane.

The reported answer and refusal cases show selected workflow behaviour, not a
reliable abstention rate. Pagination repaired selected oversized catalogue
paths, but incomplete searches and tool failures remain distinct from a
completed search without support. A late inspection also found that the global
channel can silently skip unreadable library indexes while reporting success;
until this is surfaced, a successful response is not evidence of complete
cross-library coverage. The inspection identified eight affected collections;
scoped search worked for the newly ingested example, but a probe on another
affected library returned no results. Stronger claims need
independent questions, controlled single-factor comparisons, repeated runs,
semantic support review, and tests of operations and recovery beyond the
present local setting.
